% ═══════════════════════════════════════════════════════════════════════
%  Stratégie 2 --- XAUUSD niveaux x10. Chiffres : macros générées.
% ═══════════════════════════════════════════════════════════════════════
\section{Méthodologie des tests de robustesse}
\label{app:x10_methodologie_robustesse}

Cette annexe explique, sans formalisme, ce que mesure chacun des tests de la
section~\ref{sec:x10_robustesse}, et ce qu'il ne mesure pas. Elle est écrite
pour être lisible sans bagage statistique.

\subsection{Bootstrap stationnaire par blocs}

\paragraph{La question.} Si l'histoire s'était déroulée un peu autrement,
quel résultat aurait-on obtenu~?

\paragraph{Le principe.} On ne peut pas rejouer l'histoire, mais on peut la
recoller autrement. Le bootstrap découpe la série des rendements en tronçons
de longueur aléatoire --- vingt séances en moyenne ici --- puis en tire un
grand nombre avec remise et les recolle bout à bout. On obtient une
trajectoire plausible qui n'a jamais eu lieu. Répété
\XTenBootReplicates{} fois, le procédé donne une distribution complète pour
chaque métrique, d'où l'on extrait un intervalle de confiance.

\paragraph{Pourquoi par blocs.} Un bootstrap qui tirerait les rendements un
par un détruirait la structure temporelle~: les séries de pertes, la
persistance de la volatilité, les regroupements. Les tronçons la préservent
\emph{localement}~; seul leur ordre change.

\paragraph{Ce qu'il ne dit pas.} Rien sur le futur, et rien sur le
surajustement. Il redistribue l'information de l'échantillon, il n'en crée
pas.

\subsection{Probabilistic Sharpe Ratio et Deflated Sharpe Ratio}

\paragraph{La question du PSR.} Compte tenu de la longueur de l'échantillon
et de la forme de la distribution des rendements --- asymétrie, épaisseur des
queues --- quelle est la probabilité que le vrai Sharpe soit supérieur à
zéro~?

\paragraph{La question du DSR.} La même, mais en tenant compte du nombre de
fois où l'on a cherché.

\begin{examplebox}
\textbf{Pourquoi la déflation est indispensable.} Tirez cent pièces cent fois
chacune. La meilleure série contiendra environ soixante faces, non parce que
la pièce est truquée, mais parce que vous avez pris le maximum de cent
tirages. Le même effet joue sur une grille de configurations~: le meilleur des
$N$ essais paraît d'autant meilleur que $N$ est grand.

Le \code{DSR} corrige exactement cela. Avec \XTenTrialsCount{} essais, le
Sharpe attendu du meilleur --- sous l'hypothèse qu'aucun ne vaut rien --- est
de \metric{\XTenExpectedMaxSharpe}. C'est le seuil réel, et il est positif.
\end{examplebox}

\paragraph{Le comptage des essais.} Chaque configuration évaluée est inscrite
dans un registre, avec une clé nommant l'espace de configurations exploré. Un
même balayage relancé ne compte qu'une fois~; un balayage non nommé compte
pour lui-même. C'est ce registre qui alimente la déflation --- ce qui rend le
budget d'essais opposable plutôt que déclaratif.

\paragraph{Le haircut Sharpe.} Variante du même raisonnement, appliquée à la
$p$-valeur~: elle corrige pour tests multiples et rend le Sharpe « après
correction ». Une $p$-valeur ajustée de \XTenHaircutPValue{} signifie que le
résultat est indistinguable de l'absence totale de signal.

\paragraph{La longueur minimale de backtest.} Elle répond à « combien
d'années faudrait-il pour que ce Sharpe soit crédible~? ». Elle n'est définie
que pour un Sharpe observé positif, et n'est donc pas calculable ici.

\subsection{Probability of Backtest Overfitting}

\paragraph{La question.} Si je choisis le meilleur réglage sur une moitié des
données, quelle est la probabilité qu'il se retrouve dans la moitié
\emph{inférieure} du classement sur l'autre~?

\paragraph{Le principe.} On découpe l'échantillon en un nombre pair de blocs,
on considère \emph{toutes} les manières de les répartir en deux moitiés ---
\XTenPBOSplits{} combinaisons ici --- et pour chacune, on sélectionne le
meilleur réglage sur la première moitié, puis on regarde son rang sur la
seconde. La proportion de cas où il tombe sous la médiane est le \code{PBO}.

\paragraph{La lecture.} Un \code{PBO} bas signifie que le classement des
configurations est stable~: le choix opéré sur une partie des données se
transporte. Un \code{PBO} proche de la moitié signifie que le classement est
du bruit.

\paragraph{Le piège, dans ce dossier précis.} Un \code{PBO} qui passe son
seuil sur une grille intégralement perdante n'établit aucune rentabilité ni
robustesse de la stratégie. Il dit seulement que le classement relatif
ne se retourne pas assez souvent pour franchir le seuil de surajustement. La
section~\ref{sec:x10_robustesse} le souligne, parce que c'est une confusion
facile.

\subsection{Validation croisée combinatoire purgée}

\paragraph{La question.} Que donne la stratégie sur des portions de données
qui n'ont pas servi à la choisir, en considérant \emph{plusieurs} découpages
plutôt qu'un seul~?

\paragraph{Les deux précautions.} La \emph{purge} retire les observations à
cheval sur la frontière entre apprentissage et test, pour qu'une position
ouverte d'un côté ne soit pas évaluée de l'autre. L'\emph{embargo} écarte en
plus une petite marge après la frontière, contre la persistance
d'information.

\paragraph{Ce qu'on en tire.} Non pas un chiffre, mais une distribution~: la
médiane hors échantillon de chaque configuration, ses quantiles, et la part
de découpages où elle est positive.

\subsection{Rééchantillonnage de l'ordre des transactions}

\paragraph{La question.} Le résultat vient-il de l'espérance, ou de la
séquence~?

\paragraph{Le principe.} On conserve exactement les résultats des
\XTenMCTrades{} transactions et on les remélange \XTenMCSims{} fois. Si le
repli réellement observé se situe dans la queue défavorable de la
distribution des replis simulés, la séquence a été malchanceuse. S'il se
situe au milieu, le résultat est celui qu'il devait être.

\paragraph{Le score de chance de séquence.} Il mesure cet écart en nombre
d'écarts-types. Ici il vaut \XTenSequenceZ~: l'ordre réalisé a été, si
quelque chose, marginalement \emph{plus favorable} que le tirage moyen.

\subsection{Risque de ruine}

\paragraph{La question.} Sur un grand nombre de chemins de capital plausibles,
que devient le compte~?

\paragraph{Ce qu'il faut savoir lire.} Sous un dimensionnement à fraction fixe
du capital, la probabilité de ruine \emph{littérale} est presque toujours
nulle, parce qu'une fraction d'une quantité positive reste positive. Le
chiffre informatif n'est donc pas la ruine, mais la probabilité de franchir
un seuil de perte --- ici \XTenLossFiftyProb{} de perdre plus de la moitié du
capital --- et le capital terminal médian, \XTenTerminalMedian{} du capital de
départ.

\subsection{Les seuils retenus, et leur justification}

Les seuils gelés avant la mesure sont conventionnels dans la littérature du
domaine~: au moins $0{,}95$ pour le \code{DSR}, au plus $0{,}35$ pour le
\code{PBO}, au moins 60~\% de voisins positifs pour le plateau et au moins
60~\% d'années positives pour le walk-forward.

\begin{insightbox}
\textbf{Ce qui rend un seuil valable n'est pas son niveau, c'est son
antériorité.} Un seuil de $0{,}95$ pour le \code{DSR} est discutable~; un
seuil déplacé à $0{,}80$ après avoir constaté qu'on obtenait $0{,}85$ ne l'est
pas --- il est disqualifiant. C'est pourquoi la table de décision est
committée avant la campagne, et pourquoi le hash de ce commit est reproduit
en annexe~\ref{app:x10_table_decision}.

Dans ce dossier, la question ne se pose pas~: le \code{DSR} mesuré vaut
\XTenDSR. Aucun déplacement de seuil raisonnable ne le sauverait.
\end{insightbox}

\subsection{Reproduire ces mesures}

Toutes les mesures publiées sortent d'une commande unique du dépôt, à graine
fixe, écrivant ses résultats dans des fichiers versionnés. Les tables et
figures de ce rapport sont produites par un générateur distinct qui lit
uniquement ces fichiers, et un test de la suite échoue si un chiffre publié
cesse de correspondre à sa source. La révision du dépôt ayant produit les
résultats est \code{\XTenGitHead}.
